В первой строке содержатся два целых числа $N$ и $Q$ $(1 \le N, Q \le 3 \cdot 10^5)$ --- длина массива и количество операций.

Во второй строке содержатся $Q$ целых чисел $L_1, L_2, \ldots, L_Q$ $(1 \le L_i \le N)$ --- левые границы отрезков операций.

В третьей строке содержатся $Q$ целых чисел $R_1, R_2, \ldots, R_Q$ $(L_i \le R_i \le N)$ --- правые границы отрезков операций.

В четвёртой строке содержатся $Q$ целых чисел $D_1, D_2, \ldots, D_Q$ $(|D_i| \le 10^9)$ --- числа, прибавляемые операциями: $i$-я операция прибавляет $D_i$ к каждому элементу массива с индексами от $L_i$ до $R_i$ включительно.
